\begin{statementbox}
	\textbf{\textcolor{CStatement}{条件}}
	\begin{description}[style=nextline, leftmargin=2.8em, labelsep=0.5em, itemsep=0.35em]
		\item[\textbf{\textcolor{CStatement}{条件 1（直棱柱定义）：}}] 侧棱垂直于底面，上下底面全等且平行，棱柱高为 $h$。
		\item[\textbf{\textcolor{CStatement}{条件 2（底面有外接圆）：}}] 底面多边形存在外接圆，外接圆半径为 $r$。这里默认底面至少有三个不共线顶点，且这些顶点共圆。
	\end{description}

	\textbf{\textcolor{CStatement}{结论}}
	\begin{description}[style=nextline, leftmargin=2.8em, labelsep=0.5em, itemsep=0.45em]
		\item[\textbf{\textcolor{CStatement}{结论一（球心与半径）：}}]
		      \textbf{\textcolor{CStatement}{直观描述：}} 直棱柱的外接球球心位于上、下底面外心连线的中点，半径由底面外接圆半径 $r$ 与半高 $h/2$ 通过勾股关系得到。
		      \[
			      R = \sqrt{r^{2} + \left(\frac{h}{2}\right)^{2}}.
		      \]
		      球心为上、下底面外心连线的中点。

		\item[\textbf{\textcolor{CStatement}{反求结论（已知 $R,h$ 求 $r$）：}}]
		      若已知外接球半径 $R$ 与棱柱高 $h$，则
		      \[
			      r = \sqrt{R^{2}-\left(\frac{h}{2}\right)^{2}}.
		      \]
		      使用前必须检查
		      \[
			      R\geq \frac{h}{2}.
		      \]
		      若底面为非退化多边形，则应有
		      \[
			      R>\frac{h}{2}.
		      \]

		\item[\textbf{\textcolor{CStatement}{推广结论（正棱柱特例）：}}]
		      \textbf{\textcolor{CStatement}{直观描述：}} 当底面为正 $n$ 边形，$n\geq 3$，且边长为 $a$ 时，外接圆半径可用 $a$ 与 $n$ 表示，代入即得外接球半径。
		      \[
			      r = \frac{a}{2\sin(\pi/n)},\quad
			      R = \sqrt{\frac{a^{2}}{4\sin^{2}(\pi/n)} + \frac{h^{2}}{4}}.
		      \]
	\end{description}
\end{statementbox}